% ============================================================
\chapter{PIDNet-S Model Development}
\label{ch:pidnet}
% ============================================================

\section{Selection Rationale}

The initial goal was to establish a model that could be trained, evaluated,
and plausibly deployed without requiring a high-end GPU. PIDNet-S was selected
because its architecture explicitly separates semantic context, spatial
detail, and boundary information while remaining compact
~\cite{xu2023pidnet}. This matched the main visual structure of CaT: large
connected trail regions require context, small rocks and vegetation edges
require detail, and capability boundaries are safety-relevant.

A conventional encoder--decoder such as U-Net offered a simpler alternative,
whereas a large transformer offered greater representation capacity. Neither
addressed the initial latency constraint as directly. Starting with a
purpose-built real-time network made it possible to develop the complete
pipeline around a compact deployment candidate. The comparative experiments
in Chapter~\ref{ch:rod_extension} subsequently test whether this choice
sacrifices accuracy relative to pretrained encoders and other decoder
families.

The model is trained from scratch. No Cityscapes or ImageNet checkpoint is
loaded. This keeps the baseline self-contained, but it also means later
comparisons with pretrained encoders measure complete model recipes rather
than architecture alone.

\section{Implemented Architecture}

PIDNet's control-system terminology is an analogy. The network does not
calculate a proportional--integral--derivative control command. Its three
branches instead divide the segmentation problem:
\begin{itemize}
  \item the \textbf{P branch} preserves local appearance and spatial
  resolution;
  \item the \textbf{I branch} accumulates semantic context on progressively
  smaller feature grids;
  \item the \textbf{D branch} retains high-frequency differences that help
  locate region boundaries.
\end{itemize}

After a shared stem, the P and D paths remain at one eighth of the input
resolution. The I path reaches one sixteenth, one thirty-second, and one
sixty-fourth resolution. PagFM modules compare P- and I-path features and
selectively inject semantic context into the detail path. Difference features
from the semantic stages update the D branch. PAPPM pools context at multiple
scales at the end of the I path. Light-Bag then uses boundary attention from
the D branch to fuse the spatial and semantic representations
~\cite{xu2023pidnet}.

\begin{figure}[H]
\centering
\resizebox{\textwidth}{!}{\begin{tikzpicture}[
  font=\sffamily\small,
  >=Latex,
  block/.style={draw, rounded corners=2pt, very thick, align=center,
    minimum height=1.15cm, text width=2.65cm, inner sep=5pt},
  input/.style={block, fill=gray!12},
  pblock/.style={block, draw=pidP!80!black, fill=pidP!11},
  iblock/.style={block, draw=pidI!85!black, fill=pidI!15},
  dblock/.style={block, draw=pidD!80!black, fill=pidD!12},
  fusion/.style={block, draw=pidFuse!75!black, fill=pidFuse!13},
  flow/.style={->, very thick, draw=black!70},
  semantic/.style={->, thick, draw=pidI!85!black},
  detail/.style={->, thick, dashed, draw=pidD!80!black}
]
  \node[input] (rgb) at (0,0) {\textbf{RGB input}\\$384\times640$};
  \node[input] (stem) at (3.2,0) {\textbf{Shared stem}\\features at $1/8$};

  \node[pblock] (p) at (6.5,3.0) {\textbf{P branch}\\spatial detail at $1/8$\\two PagFM updates};
  \node[iblock] (i) at (6.5,0) {\textbf{I branch}\\context from $1/16$ to $1/64$\\PAPPM pooling};
  \node[dblock] (d) at (6.5,-3.0) {\textbf{D branch}\\boundary cues at $1/8$\\semantic differences};

  \node[fusion] (bag) at (10.2,0) {\textbf{Light-Bag}\\boundary-guided\\P/I fusion};
  \node[fusion] (head) at (13.6,0) {\textbf{Segmentation head}\\two-class logits\\at $1/8$};

  \draw[flow] (rgb) -- (stem);
  \draw[flow] (stem) -- (p);
  \draw[flow] (stem) -- (i);
  \draw[flow] (stem) -- (d);

  \draw[flow] (p) -- (bag);
  \draw[flow] (i) -- (bag);
  \draw[flow] (d) -- (bag);
  \draw[flow] (bag) -- (head);
\end{tikzpicture}
}
\caption[Conceptual PIDNet-S three-branch architecture.]{Conceptual PIDNet-S
flow. The P branch retains spatial detail, the I branch extracts context, and
the D branch supplies boundary cues. PagFM, PAPPM, and Light-Bag exchange and
fuse those representations before the segmentation head. Original schematic
drawn for this thesis and adapted from Xu et al.~\cite{xu2023pidnet}.}
\label{fig:final_pidnet_architecture}
\end{figure}

Figure~\ref{fig:final_pidnet_architecture} shows the three processing branches
and the points at which they exchange and fuse features.

The implementation uses the PIDNet-S schedule $(m,n,C)=(2,3,32)$, with 96
pyramid-pooling channels and 128 head channels. $C$ is the base width; $m$ and
$n$ control residual-stage repetition. Table~\ref{tab:pidnet_tensor_shapes}
records the principal shapes for the $640\times384$ input. Listing these
shapes is useful both for implementation verification and for understanding
where computation is spent.

\begin{table}[H]
\centering
\caption{Principal implemented PIDNet-S feature shapes at $640\times384$.}
\label{tab:pidnet_tensor_shapes}
\small
\begin{tabularx}{\textwidth}{@{}lclX@{}}
\toprule
\textbf{Component} & \textbf{Spatial size} & \textbf{Channels} & \textbf{Purpose} \\
\midrule
Shared stem output & $80\times48$ ($1/8$) & 64 & Common input to the three paths \\
P branch & $80\times48$ ($1/8$) & $64\rightarrow128$ & Preserve detail; receive two PagFM updates \\
I stage 3 & $40\times24$ ($1/16$) & 128 & First deeper semantic stage \\
I stage 4 & $20\times12$ ($1/32$) & 256 & Broader contextual representation \\
I stage 5 & $10\times6$ ($1/64$) & 512 & Deepest semantic features before PAPPM \\
PAPPM output & $80\times48$ ($1/8$) & 128 & Multi-scale context returned for fusion \\
D branch & $80\times48$ ($1/8$) & $32\rightarrow64\rightarrow128$ & Boundary-sensitive representation \\
Light-Bag output & $80\times48$ ($1/8$) & 128 & Boundary-guided P/I fusion \\
Binary head & $80\times48$ ($1/8$) & 2 & Native untraversable/traversable logits \\
Wrapper output & $640\times384$ & 2 & Bilinearly upsampled full-resolution logits \\
\bottomrule
\end{tabularx}
\end{table}

Auxiliary P- and D-head outputs are disabled, so the network exposes one final
two-class logits tensor and has 7,623,522 trainable parameters. In the vendored
PIDNet implementation this architectural choice is named
\texttt{augment=False}. It is unrelated to the data-loader
\texttt{augment} flag. Both happen to be false for the retained checkpoint, but
they control different operations.

\begin{figure}[H]
\centering
\resizebox{\textwidth}{!}{\begin{tikzpicture}[
  font=\sffamily\small,
  >=Latex,
  node distance=0.7cm,
  stage/.style={draw, rounded corners=2pt, very thick, align=center,
    minimum height=1.2cm, text width=2.8cm, inner sep=5pt},
  data/.style={stage, fill=gray!12},
  model/.style={stage, draw=pidP!80!black, fill=pidP!11},
  operation/.style={stage, draw=pidI!85!black, fill=pidI!14},
  output/.style={stage, draw=pidFuse!75!black, fill=pidFuse!13},
  target/.style={stage, draw=pidD!80!black, fill=pidD!11},
  flow/.style={->, very thick, draw=black!70},
  trainflow/.style={->, thick, dashed, draw=pidD!80!black}
]
  \node[data] (rgb) {RGB tensor\\$3\times384\times640$};
  \node[model, right=of rgb] (core) {\textbf{PIDNet-S core}\\P + I + D branches};
  \node[operation, right=of core] (logits) {\textbf{Two-class logits}\\$2\times48\times80$\\bilinear resize to input};
  \node[output, right=of logits] (mask) {\textbf{Deployment decision}\\softmax; $p_{\mathrm{trav}}\geq0.50$\\binary mask};

  \draw[flow] (rgb) -- (core);
  \draw[flow] (core) -- (logits);
  \draw[flow] (logits) -- (mask);

  \node[target, below=1.25cm of logits] (labels) {\textbf{CaT target}\\IDs 1,2 $\rightarrow$ traversable\\IDs 0,3 $\rightarrow$ untraversable};
  \node[target, below=1.25cm of core] (loss) {\textbf{Training only}\\weighted cross-entropy\\+ Dice loss};
  \draw[trainflow] (labels) -- (loss);
  \draw[trainflow] (loss) -- node[left, align=right] {updates\\weights} (core);
\end{tikzpicture}
}
\caption[Implemented PIDNet-S binary pipeline.]{Implemented PIDNet-S binary
pipeline. The one-eighth-resolution logits are returned to the input
resolution before the training loss or deployment decision is calculated.
The lower dashed path is active only during training. Original schematic drawn
for this thesis; the core network follows Xu et al.~\cite{xu2023pidnet}.}
\label{fig:final_pidnet_binary_pipeline}
\end{figure}

Figure~\ref{fig:final_pidnet_binary_pipeline} then connects that core network
to the two-class training and thresholding operations implemented here.

For logits $z_{i0}$ and $z_{i1}$ at pixel $i$, the traversable probability is
\begin{equation}
 p_i =
 \frac{\exp(z_{i1})}
 {\exp(z_{i0})+\exp(z_{i1})}.
 \label{eq:pidnet_softmax}
\end{equation}
Here, $i$ indexes a pixel, $z_{i0}$ and $z_{i1}$ are its untraversable and
traversable logits, $\exp(\cdot)$ is the exponential function, and $p_i$ is
the normalised probability assigned to the traversable class. The other
softmax output, $1-p_i$, is the untraversable probability.
The reference binary prediction is $\hat y_i=\mathbb{1}[p_i\geq0.50]$.
The constant 0.50 is the fixed probability threshold, and equality is assigned
to the traversable class. No morphological cleanup, connected-component filtering, or temporal smoothing
is applied. The output therefore represents the model and threshold directly.

\section{Objective, Optimiser, and Checkpoint Selection}

The main PIDNet objective combines weighted cross-entropy and soft Dice loss.
For classes $c\in\{0,1\}$ and their corresponding weights $w_c$, the
cross-entropy term is
\begin{equation}
 \mathcal{L}_{\mathrm{WCE}}
 =-\frac{1}{N}\sum_{i=1}^{N}w_{y_i}\log p(y_i\mid x).
\end{equation}
In this expression, $\mathcal{L}_{\mathrm{WCE}}$ is weighted cross-entropy,
$N$ is the number of pixels included in the loss, $i$ indexes those pixels,
$y_i\in\{0,1\}$ is the target class, $w_{y_i}$ is the weight of that class,
and $p(y_i\mid x)$ is the probability assigned to the correct class for input
$x$. The natural logarithm is denoted by $\log$.
The implemented Dice term operates on the traversable probability and binary
target. For one image,
\begin{equation}
 \mathcal{L}_{\mathrm{Dice}}
 =1-\frac{2\sum_i p_i y_i+\epsilon}
 {\sum_i p_i+\sum_i y_i+\epsilon},
 \label{eq:pidnet_dice}
\end{equation}
Here, $p_i$ is the traversable probability from
Equation~\ref{eq:pidnet_softmax}, $y_i$ is one for a traversable target and
zero otherwise, each sum covers the pixels of one image, and
$\epsilon=10^{-6}$ prevents division by zero. The resulting
$\mathcal{L}_{\mathrm{Dice}}$ is zero for perfect overlap and approaches one
as overlap is lost. The constant factor 2 is the conventional Dice
normalisation: when prediction and target coincide, the doubled intersection
equals the sum of their two sizes. The batch mean is used. Dice complements per-pixel
cross-entropy with a
region-overlap signal~\cite{milletari2016vnet}. The default total is
\begin{equation}
 \mathcal{L}=\mathcal{L}_{\mathrm{WCE}}+\mathcal{L}_{\mathrm{Dice}}.
\end{equation}
Thus, $\mathcal{L}$ is the scalar objective differentiated during the default
training run, and both component losses enter with unit coefficient.

Class weights are computed from inverse training-pixel frequency. Before
normalisation, the untraversable weight is multiplied by 2.5 and the
traversable weight by 0.75. The two final weights are rescaled to sum to two.
For the fixed split they are $[1.2606075,\,0.7393925]$. The intent was to
discourage false-safe predictions by making errors on the blocked class more
expensive. Because loss weighting changes optimisation rather than directly
optimising FSR, its actual effect must be established experimentally.

An optional false-safe probability penalty was also implemented:
\begin{equation}
 \mathcal{L}_{\mathrm{FS}}
 =\frac{\sum_i p_i\,\mathbb{1}[y_i=0]}
 {\sum_i\mathbb{1}[y_i=0]}.
 \label{eq:false_safe_penalty}
\end{equation}
$i$ indexes the same $N$ loss pixels as above, and both sums run over those
pixels. The indicator $\mathbb{1}[y_i=0]$ selects annotated untraversable pixels, so
the numerator sums their predicted traversable probabilities and the
denominator counts them. Consequently, $\mathcal{L}_{\mathrm{FS}}$ is the
mean traversable probability over the blocked class. With non-negative
coefficient $\lambda_{\mathrm{FS}}$, the tested objective becomes
$\mathcal{L}+\lambda_{\mathrm{FS}}\mathcal{L}_{\mathrm{FS}}$; the ablation uses
$\lambda_{\mathrm{FS}}=0.35$.
This term pushes traversable probability down over all annotated negative
pixels. It was a plausible alternative to class weighting, but the controlled
experiment later showed that the tested coefficient did not improve the
desired rate.

The retained training settings are given in
Table~\ref{tab:frozen_pidnet_training}. AdamW supplies decoupled weight decay
~\cite{loshchilov2019adamw}. Cosine annealing reduces the learning rate from
$2.5\times10^{-4}$ toward $10^{-5}$ over 15 epochs. Bfloat16 automatic mixed
precision is used on CUDA when supported, while model parameters and optimiser
state remain stored in FP32.

\begin{table}[H]
\centering
\caption{Training configuration of the reference PIDNet-S model.}
\label{tab:frozen_pidnet_training}
\small
\begin{tabularx}{\textwidth}{@{}lX@{}}
\toprule
\textbf{Setting} & \textbf{Implemented value and function} \\
\midrule
Initialisation & Random PIDNet-S weights; no pretrained checkpoint \\
Input and batch & $640\times384$ RGB; batch size 4 \\
Epochs & 15 complete passes over 1,002 training images \\
Optimiser & AdamW, initial learning rate $2.5\times10^{-4}$, weight decay $10^{-4}$ \\
Scheduler & Cosine annealing by epoch to minimum learning rate $10^{-5}$ \\
Objective & weighted cross-entropy + traversable soft Dice \\
Class weights & $[1.2606075,\,0.7393925]$ for untraversable/traversable \\
Data augmentation & disabled \\
Mixed precision & bfloat16 autocast on supported CUDA hardware \\
Seed and workers & seed 1337; zero loader worker processes in the portable config \\
Selection & save the epoch with maximum validation mIoU \\
Decision threshold & 0.50, fixed before official test evaluation \\
\bottomrule
\end{tabularx}
\end{table}

After each epoch, all 266 validation images are evaluated without augmentation
or shuffling. A new \texttt{best.pt} is written only if validation mIoU
increases. The retained checkpoint comes from epoch 15, where validation mIoU
is 0.88834. This rule prevents the 544-image test score from participating in
model selection.

\section{Exploratory Development}

The earliest experiments were iterative engineering runs. They are useful
because they show how the system changed, but they are not single-factor
ablations: several settings change together. Table
~\ref{tab:pidnet_iterative_development} lists both configuration and test
outcome so that the later controlled study is not confused with this
exploration.

\begin{table}[H]
\centering
\caption{Exploratory PIDNet-S development runs. Values are test results; the
runs are confounded and do not estimate individual causal effects.}
\label{tab:pidnet_iterative_development}
\footnotesize
\setlength{\tabcolsep}{4pt}
\begin{tabularx}{\textwidth}{@{}p{0.17\textwidth}Xlrrrr@{}}
\toprule
\textbf{Run} & \textbf{Main change(s)} & \textbf{Policy} & \textbf{mIoU} & \textbf{FSR} & \textbf{FBR} & \textbf{Thr.} \\
\midrule
\texttt{default} & Initial pipeline, five epochs & Blue & 0.8374 & 7.24\% & 7.92\% & 0.50 \\
\texttt{conservative} & 15 epochs and biased class weights & Blue & 0.8735 & 3.89\% & 9.67\% & 0.50 \\
\texttt{third\_run} & $720\times448$, false-safe term, 18 epochs & Blue & 0.8219 & 2.32\% & 21.42\% & 0.58 \\
\texttt{fourth\_run} & Remove false-safe term; retain higher resolution & Blue & 0.8544 & 4.08\% & 12.30\% & 0.53 \\
\path{blue_green_best} & Change policy; $640\times384$, 15 epochs & Blue+green & 0.8780 & 4.87\% & 8.34\% & 0.50 \\
\bottomrule
\end{tabularx}
\end{table}

The five-epoch run established that the complete pipeline functioned but left
clear room for learning. Extending training and biasing the loss reduced FSR
while increasing FBR, suggesting a more conservative operating point. The
third run combined higher resolution, an explicit false-safe penalty, and a
0.58 threshold. Its low FSR was accompanied by a very high 21.42\% FBR and
lower mIoU. Because three factors changed, it was impossible to identify the
cause. The fourth run removed the explicit penalty but retained the higher
resolution and a threshold above 0.50; it recovered part of the accuracy but
remained overly conservative.

The blue+green run changed the target vehicle and returned to
$640\times384$. Its result motivated using blue+green as the main operational
policy, but it does not prove that merging green causes an accuracy
improvement: target definition, training history, and class balance differ
from the earlier blue-only runs. This interpretive constraint is retained
throughout the thesis.

The exploratory phase produced two lessons. First, a low FSR can be obtained
by simply blocking more pixels, so FSR must always be read with FBR. Second,
multi-factor iteration can guide engineering but cannot support a claim about
which setting caused the change. Those lessons led to the controlled ablation.

\section{Controlled Configuration Ablation}

The controlled reference uses the blue+green policy, PIDNet-S,
$640\times384$, batch size 4, 15 epochs, seed 1337, augmentation enabled,
biased class weights, CE+Dice, and threshold 0.50. Four variants each change
one factor:
\begin{enumerate}
  \item add the false-safe penalty with weight 0.35;
  \item disable the complete tested augmentation recipe;
  \item increase input to $720\times448$ and reduce batch size to 2 for memory;
  \item replace the computed class weights with neutral $[1,1]$.
\end{enumerate}

The alternatives were chosen to answer practical questions arising from the
exploratory work. The penalty tests whether FSR can be reduced inside the loss
rather than at the threshold. Disabling augmentation tests whether the
synthetic appearance variation matches CaT. Higher resolution tests whether
thin boundaries are limited by input detail. Neutral weights test whether the
manually biased imbalance correction is necessary.

\begin{table}[H]
\centering
\caption{Single-seed controlled PIDNet-S ablation. Bold identifies the
configuration retained for subsequent analysis. Lower FSR and FBR are
preferable.}
\label{tab:pidnet_controlled_ablation}
\small
\begin{tabularx}{\textwidth}{@{}lXrrr@{}}
\toprule
\textbf{Run} & \textbf{One changed factor} & \textbf{mIoU} & \textbf{FSR} & \textbf{FBR} \\
\midrule
Reference & None & 0.8667 & 5.89\% & 8.48\% \\
False-safe loss & Penalty coefficient 0.35 & 0.8177 & 7.45\% & 12.95\% \\
\textbf{Augmentation off} & \textbf{Disable tested flip and photometric recipe} & \textbf{0.8939} & \textbf{4.32\%} & 7.04\% \\
Higher resolution & $720\times448$, batch size 2 & 0.8632 & 5.62\% & 9.28\% \\
Neutral weights & Class weights $[1,1]$ & 0.8848 & 5.28\% & \textbf{6.95\%} \\
\bottomrule
\end{tabularx}
\end{table}

Table~\ref{tab:pidnet_controlled_ablation} supplies the test result for each
one-factor change and marks the retained configuration.

Disabling the tested augmentation improves mIoU by 2.72 percentage points and
reduces both asymmetric error rates relative to the reference. The result
suggests that the colour, blur, and sharpness variations used here do not
represent helpful additional CaT variation. It cannot establish that a
different geometric, weather, or sensor-specific augmentation would also be
harmful.

The false-safe penalty is discarded because it worsens mIoU, FSR, and FBR.
Its gradient pushes traversable probability downward on every negative pixel,
but that global pressure does not guarantee a better thresholded operating
point after optimisation. Higher resolution is also discarded: it increases
pixel computation by 31.25\%, halves the batch size, and does not improve any
headline metric. Neutral weights produce the lowest FBR but a higher FSR and
lower mIoU than augmentation-off. The augmentation-off checkpoint is selected
by validation mIoU. Settings from different rows are not combined after
examining their test results.

Because each row has one seed, the differences provide engineering evidence
rather than estimates of population effects. The separate three-seed
architecture experiment measures repeatability for complete model recipes; it
does not change the statistical scope of this ablation.

\section{Threshold as an Operating Control}

A threshold sweep was conducted on the earlier
\texttt{blue\_green\_best} checkpoint from 0.35 to 0.65. It was not used to
alter the reference checkpoint's 0.50 threshold. Table
~\ref{tab:pidnet_threshold_sweep} shows the expected monotonic trade:
increasing the threshold reduces false-safe openings and increases false
blocks.

\begin{table}[H]
\centering
\caption{Post-hoc threshold sweep on the earlier
\texttt{blue\_green\_best} checkpoint.}
\label{tab:pidnet_threshold_sweep}
\begin{tabular}{@{}rrrr@{}}
\toprule
\textbf{Threshold} & \textbf{mIoU} & \textbf{FSR} & \textbf{FBR} \\
\midrule
0.35 & 0.8774 & 6.14\% & 6.78\% \\
0.40 & 0.8781 & 5.68\% & 7.29\% \\
0.45 & 0.8782 & 5.26\% & 7.81\% \\
0.50 & 0.8780 & 4.87\% & 8.34\% \\
0.55 & 0.8774 & 4.50\% & 8.89\% \\
0.60 & 0.8763 & 4.15\% & 9.49\% \\
0.65 & 0.8747 & 3.80\% & 10.14\% \\
\bottomrule
\end{tabular}
\end{table}

\begin{figure}[H]
\centering
\begin{tikzpicture}
\begin{axis}[
  width=0.84\textwidth,
  height=6.5cm,
  xlabel={Traversable probability threshold},
  ylabel={Pixel error rate},
  ymin=0.03,ymax=0.105,
  xmin=0.34,xmax=0.66,
  xtick={0.35,0.40,0.45,0.50,0.55,0.60,0.65},
  xticklabels={0.35,0.40,0.45,0.50,0.55,0.60,0.65},
  ytick={0.04,0.06,0.08,0.10},
  yticklabels={4\%,6\%,8\%,10\%},
  grid=both,
  grid style={line},
  legend style={
    at={(0.5,1.04)},
    anchor=south,
    legend columns=2,
    fill=white,
    draw=black!35
  }
]
\addplot[color=safety,mark=*,thick] coordinates {
  (0.35,0.0614) (0.40,0.0568) (0.45,0.0526) (0.50,0.0487)
  (0.55,0.0450) (0.60,0.0415) (0.65,0.0380)
};
\addlegendentry{FSR}
\addplot[color=block,mark=square*,thick] coordinates {
  (0.35,0.0678) (0.40,0.0729) (0.45,0.0781) (0.50,0.0834)
  (0.55,0.0889) (0.60,0.0949) (0.65,0.1014)
};
\addlegendentry{FBR}
\end{axis}
\end{tikzpicture}
\caption{False-safe and false-block trade-off under post-hoc threshold
variation. A threshold is an operating choice; it does not improve both error
directions simultaneously.}
\label{fig:pidnet_threshold_plot}
\end{figure}

Figure~\ref{fig:pidnet_threshold_plot} makes the opposing movement of FSR and
FBR across the tested thresholds directly visible.

The maximum observed mIoU occurs at 0.45, but the difference from 0.50 is only
0.0002. More importantly, the sweep uses the test set for characterisation.
Choosing 0.45 from this table and reporting it as an untouched test result
would leak test information into model selection. The final 0.50 operating
point is retained. A deployed threshold should instead be selected on
validation data using an application-defined error cost or risk constraint, then
evaluated once on an independent test set.

\section{Reference Checkpoint}

All subsequent PIDNet-S error analysis, reproduction, CPU, and TensorRT measurements
use the following checkpoint:
\begin{center}
\small\path{outputs/controlled_ablation_augment_off/checkpoints/best.pt}
\end{center}
Keeping the parameters fixed preserves a direct link between the configuration,
manifest, decision threshold, pixel metrics, qualitative panels, and exported
engines. The checkpoint is the strongest PIDNet-S configuration observed in
the controlled study, not an estimate of the architecture's maximum possible
performance. Other augmentation families, pretrained weights, boundary
supervision, longer schedules, or different seed sets could produce a
different result. Its value is that it identifies one model whose behaviour
can be reproduced across evaluation and deployment settings.
